\begin{pastexam}[label={px:ust-107-ee-ds-06}]{臺聯大電機類 107 學年度 第六題}
（問答；10 分）

(10\%) Given a list $L$ of length $n$.
\begin{enumerate}
\item[(一)] (2\%) What is the time complexity of searching an element in $L$ using linear search?
\item[(二)] (2\%) What is the time complexity of searching an element in sorted $L$ using binary search?
\item[(三)] (6\%) When does it make sense to perform a comparison-based sorting on $L$ first then binary searches? Assume we do $k$ searches.
\end{enumerate}
\end{pastexam}

\begin{pastsoln}
(一)~$O(n)$。 線性搜尋從頭比到找到為止；最壞（不在 $L$ 裡或在最後一個）要比 $n$ 次，找得到時平均約 $n/2$ 次，
都是 $\Theta(n)$。

(二)~$O(\log n)$。 二分搜尋每比一次，剩下的範圍減半，最多比 $\lfloor \lg n \rfloor + 1$ 次（資結書第~17.1~節）。
前提是 $L$ 已經排好，而且能用索引直接存取第 $i$ 個元素（陣列）；若 $L$ 是鏈結串列，找中間那個就要 $\Theta(n)$，
二分搜尋沒有好處。

(三)~比較兩種做法 $k$ 次搜尋的總時間：
\begin{itemize}
\item 不排序、每次線性搜尋：$k \cdot \Theta(n) = \Theta(kn)$。
\item 先排序再二分搜尋：比較排序至少 $\Omega(n \log n)$（資結書定理~15.1），用合併排序或堆積排序就是
$\Theta(n \log n)$；之後每次 $O(\log n)$。 總共 $\Theta(n \log n + k \log n)$。
\end{itemize}
先排序划算，要
\[
n \log n + k \log n < c\, kn
\quad\Longleftrightarrow\quad
k > \frac{n \log n}{c\,n - \log n} \approx \frac{\log n}{c} ,
\]
其中 $c$ 是兩邊常數的比例。 所以結論是：\textbf{搜尋次數 $k$ 比 $\log n$ 多（$k = \Omega(\log n)$，更精確地說 $k$ 超過
某個常數乘 $\log n$）時，先排序再二分搜尋比較快}；$k$ 很小（例如只查一兩次、或 $k = o(\log n)$）時，排序的
$n \log n$ 攤不掉，直接線性搜尋比較好。 $k$ 很大時差距很明顯：$k = n$ 次搜尋，線性是 $\Theta(n^2)$，先排序是
$\Theta(n \log n)$。

例如 $n = 10^6$，$\lg n \approx 20$：排序大約 $2 \times 10^7$ 次比較，一次線性搜尋平均 $5 \times 10^5$ 次，
所以要查 $40$ 次以上，排序才開始划算。

另外兩個實務條件：$L$ 在多次搜尋之間不會一直變動（每插入一個元素就要維持排序，或改用平衡搜尋樹）；以及可以
改變 $L$ 的順序（或有空間複製一份來排）。
\end{pastsoln}
